% Lecture 9: determinant applications and official midterm solutions.
% Compile from the repository root with:
%   pdflatex -interaction=nonstopmode -halt-on-error -output-directory=Lecture9 Lecture9/lecture9.tex
\pdfminorversion=7
\documentclass[aspectratio=169]{beamer}

\usepackage{amsmath,amssymb,mathtools}
\usepackage{array,booktabs}
\usepackage{xcolor}
\usepackage{iftex}
\ifPDFTeX
  \usepackage[UTF8,heading=false]{ctex}
\else
  \usepackage{fontspec}
  \usepackage{xeCJK}
  \setCJKmainfont{Microsoft YaHei}
\fi

\makeatletter
\def\input@path{{./}{../}{Lecture9/}}
\makeatother

\mode<presentation>{
  \usetheme{OSUCSE169}
  \setbeamercovered{invisible}
}
\setbeamertemplate{navigation symbols}{}
\setbeamertemplate{footline}[frame number]
\setbeamercolor{headline}{fg=black,bg=black}
\setbeamercolor{section in head/foot}{bg=gray}
\setbeamertemplate{headline}{%
  \begin{beamercolorbox}[ht=2.25ex,dp=3.75ex]{section in head/foot}
    \insertnavigation{\paperwidth}
  \end{beamercolorbox}}

\definecolor{lecturegreen}{RGB}{48,105,21}
\definecolor{lectureblue}{RGB}{0,76,170}
\definecolor{lectureorange}{RGB}{202,105,0}

\newcommand{\R}{\mathbb{R}}
\newcommand{\vect}[1]{\mathbf{#1}}
\newcommand{\rank}{\operatorname{rank}}
\newcommand{\sectioncard}[2]{%
  \begin{frame}[plain]
    \vfill\centering
    {\Huge\color{lectureorange}\bfseries #1}\par\bigskip
    {\LARGE\color{lecturegreen}\bfseries #2}
    \vfill
  \end{frame}}
\newcommand{\theorembox}[2]{%
  \begin{beamercolorbox}[sep=1.1ex,wd=\linewidth]{block body}
    \textbf{#1}\par #2
  \end{beamercolorbox}}
\newcommand{\answerbox}[1]{%
  \begin{beamercolorbox}[sep=1.0ex,wd=\linewidth]{block body}
    \textbf{Answer:} #1
  \end{beamercolorbox}}

\title{MATH 100: Linear Algebra}
\subtitle{Lecture 9 - Determinant Applications and Midterm Solutions}
\author{Yuqing Yang}
\institute{Macau University of Science and Technology}
\date{10/30/2026}

\begin{document}

\begin{frame}[plain]
  \titlepage
\end{frame}

% \begin{frame}{Today's session}
% \begin{enumerate}
%   \item Complete Section 3.3: applications of determinant properties.
%   \item Review the official midterm answers and the shortest reliable methods.
%   \item Identify the steps that earn partial credit on worked problems.
% \end{enumerate}
% \end{frame}

% \begin{frame}{Learning goals}
% By the end of the lecture, you should be able to
% \begin{itemize}
%   \item connect determinants with invertibility and the number of solutions;
%   \item use determinant properties to simplify computations;
%   \item apply Cramer's rule to small square systems;
%   \item explain each official midterm answer with enough supporting work.
% \end{itemize}
% \end{frame}

% -----------------------------------------------------------------------------
\section[Determinant Applications]{Applications of Determinant Properties}
\sectioncard{3.3}{Applications of Determinant Properties}

\begin{frame}{Four applications}
\begin{enumerate}
  \item Efficient determinant computation
  \item The invertibility test
  \item Classification of square linear systems
  \item Cramer's rule
\end{enumerate}
\medskip
The determinant gives both a computational tool and a structural test for a square matrix.
\end{frame}

\begin{frame}{Efficient determinant computation}
Two ideas work well together:
\begin{enumerate}
  \item create zeros with row or column replacement;
  \item use cofactor expansion or finish with a triangular determinant.
\end{enumerate}
\pause
Before computing, check for
\begin{itemize}
  \item a row or column with many zeros;
  \item a quick replacement that creates zeros;
  \item repeated or proportional rows or columns.
\end{itemize}
\end{frame}

\begin{frame}{Efficient computation example}
\small
Compute
\[
D=\begin{vmatrix}
1&4&5&0\\4&5&0&1\\5&0&1&4\\0&1&4&5
\end{vmatrix}.
\]
Adding all four rows to the first row gives $(10,10,10,10)$, so
\[
D=10\begin{vmatrix}
1&1&1&1\\4&5&0&1\\5&0&1&4\\0&1&4&5
\end{vmatrix}.
\]
Row replacement and cofactor expansion reduce the expression to
\[
10\cdot8\cdot(-8)\begin{vmatrix}3&2\\1&1\end{vmatrix}
=\boxed{-640}.
\]
\end{frame}

\begin{frame}{The invertibility test}
\theorembox{Invertibility theorem}{A square matrix $A$ is invertible if and only if $\det A\ne0$.}
\medskip
If $A$ is invertible, then
\[
1=\det I=\det(AA^{-1})=(\det A)(\det A^{-1}),
\]
so $\det A\ne0$.
\pause

Elimination shows that $\det A\ne0$ forces a pivot in every row and column. Therefore $A$ is invertible.
\end{frame}

\begin{frame}{The determinant of an inverse}
If $A$ is invertible, then
\[
(\det A)(\det A^{-1})=\det I=1.
\]
Therefore
\[
\boxed{\det(A^{-1})=\frac{1}{\det A}}.
\]
\pause
If $AB=I$ for square matrices $A$ and $B$, then $\det A\det B=1$. Both matrices are invertible and $B=A^{-1}$.
\end{frame}

\begin{frame}{A singular product}
Suppose $A$ and $B$ are $n\times n$ matrices, $\det B\ne0$, and $\det(AB)=0$.
Then
\[
0=\det(AB)=(\det A)(\det B).
\]
Because $\det B\ne0$,
\[
\boxed{\det A=0}.
\]
Thus $A$ is singular.
\end{frame}

\begin{frame}[shrink=4]{Equivalent conditions for a square matrix}
For an $n\times n$ matrix $A$, the following statements are equivalent:
\begin{enumerate}
  \item $A$ is invertible.
  \item $A$ is row equivalent to $I_n$.
  \item $\det A\ne0$.
  \item $A\vect x=\vect b$ has a unique solution for every $\vect b\in\R^n$.
  \item $A\vect x=\vect0$ has only the trivial solution.
  \item $A^T$ is invertible.
  \item $\rank A=n$.
\end{enumerate}
\end{frame}

\begin{frame}{Classifying square systems}
For a square system $A\vect x=\vect b$:
\begin{itemize}
  \item if $\det A\ne0$, the system has exactly one solution;
  \item if $\det A=0$, compare $\rank A$ and $\rank[A\mid\vect b]$.
\end{itemize}
\pause
When $\det A=0$,
\[
\begin{array}{ccl}
\rank A=\rank[A\mid\vect b]<n &\Longrightarrow& \text{infinitely many solutions},\\[1ex]
\rank A<\rank[A\mid\vect b] &\Longrightarrow& \text{no solution}.
\end{array}
\]
\end{frame}

\begin{frame}{A parameterized system}
Classify the solutions of
\[
\left\{\begin{aligned}
kx+y+z&=1,\\
x+ky+z&=k,\\
x+y+kz&=k^2.
\end{aligned}\right.
\]
The coefficient matrix is
\[
A=\begin{bmatrix}k&1&1\\1&k&1\\1&1&k\end{bmatrix},
\qquad
\det A=(k+2)(k-1)^2.
\]
\end{frame}

\begin{frame}{Parameterized-system classification}
From $\det A=(k+2)(k-1)^2$:
\begin{itemize}
  \item If $k\ne1,-2$, the system has a unique solution.
  \item If $k=1$, then $\rank A=\rank[A\mid\vect b]=1<3$, so the system has infinitely many solutions.
  \item If $k=-2$, then $\rank A=2<\rank[A\mid\vect b]=3$, so the system has no solution.
\end{itemize}
\end{frame}

\begin{frame}{Cramer's rule}
\theorembox{Cramer's rule}{Let $A$ be an invertible $n\times n$ matrix. The unique solution of $A\vect x=\vect b$ is}
\[
\boxed{x_i=\frac{\det A_i(\vect b)}{\det A}},
\qquad i=1,2,\ldots,n,
\]
where $A_i(\vect b)$ replaces column $i$ of $A$ with $\vect b$.
\pause
Cramer's rule requires $\det A\ne0$.
\end{frame}

\begin{frame}{Cramer's rule example}
Solve
\[
\left\{\begin{aligned}
x_1+2x_2-x_3&=-5,\\
2x_1-x_2+x_3&=6,\\
x_1-x_2-3x_3&=-3.
\end{aligned}\right.
\]
The determinants are
\[
\det A=19,
\quad \det A_1=19,
\quad \det A_2=-38,
\quad \det A_3=38.
\]
Thus
\[
\boxed{(x_1,x_2,x_3)=(1,-2,2)}.
\]
\end{frame}

\begin{frame}{A four-variable example}
\small
For the system
\[
\left\{\begin{aligned}
2x_1+x_2-5x_3+x_4&=8,\\
x_1-3x_2-6x_4&=9,\\
2x_2-x_3+2x_4&=-5,\\
x_1+4x_2-7x_3+6x_4&=0,
\end{aligned}\right.
\]
the determinants are
\[
\det A=27,
\qquad
(\det A_1,\det A_2,\det A_3,\det A_4)=(81,-108,-27,27).
\]
Therefore
\[
\boxed{(x_1,x_2,x_3,x_4)=(3,-4,-1,1)}.
\]
\end{frame}

\begin{frame}{When Cramer's rule is useful}
Cramer's rule gives an exact formula and supports theoretical proofs.
\medskip

For hand calculations, it is practical mainly for $2\times2$ and small $3\times3$ systems. Larger systems require many determinants, so elimination is usually more efficient.
\medskip

Use row reduction for a general linear system. Use inverse matrices or Cramer's rule only when the coefficient matrix is square and invertible.
\end{frame}

% -----------------------------------------------------------------------------
\section[Midterm Solutions]{Official Midterm Solutions}
% \sectioncard{Midterm}{Official Answers and Walkthrough}

% \begin{frame}{Walkthrough format}
% For each question, focus on
% \begin{itemize}
%   \item the idea being tested;
%   \item the shortest reliable method;
%   \item the work needed to justify the answer;
%   \item the final answer in a form that can be checked quickly.
% \end{itemize}
% \end{frame}

\subsection{Multiple choice}
\sectioncard{Section I}{Multiple Choice - 20 points}

\begin{frame}{MC 1 - Row echelon form}
In row echelon form,
\begin{itemize}
  \item all zero rows appear at the bottom;
  \item each leading entry lies to the right of the leading entry above it;
  \item entries below each pivot are zero.
\end{itemize}
Only
\[
\begin{bmatrix}1&2&0\\0&0&1\\0&0&0\end{bmatrix}
\]
satisfies all three conditions.
\answerbox{\textbf{B}.}
\end{frame}

\begin{frame}{MC 2 - Row replacement}
To change
\[
\begin{bmatrix}1&2\\3&7\end{bmatrix}
\quad\text{into}\quad
\begin{bmatrix}1&2\\0&1\end{bmatrix},
\]
replace the second row:
\[
(3,7)-3(1,2)=(0,1).
\]
\answerbox{\textbf{B}. $R_2\leftarrow R_2-3R_1$.}
\end{frame}

\begin{frame}{MC 3 - Evaluating a transformation}
Given
\[
T(x,y)=(2x-y,\ x+3y),
\]
substitute $(x,y)=(1,-2)$:
\[
T(1,-2)=(2+2,\ 1-6)=(4,-5).
\]
\answerbox{\textbf{B}. $(4,-5)$.}
\end{frame}

\begin{frame}{MC 4 - Matrix dimensions}
Suppose $A$ is $2\times3$, $B$ is $3\times2$, and $C$ is $2\times2$.
\begin{align*}
AB&:(2\times3)(3\times2)\text{ is defined},\\
BA&:(3\times2)(2\times3)\text{ is defined},\\
CA&:(2\times2)(2\times3)\text{ is defined}.
\end{align*}
Addition requires equal dimensions, but $A$ and $B$ have different sizes.
\answerbox{\textbf{A}. $A+B$ is not defined.}
\end{frame}

\begin{frame}{MC 5 - Isolating a matrix factor}
Starting with
\[
XA=C,
\]
multiply on the right by $A^{-1}$:
\[
XAA^{-1}=CA^{-1}.
\]
Since $AA^{-1}=I$,
\[
X=CA^{-1}.
\]
\answerbox{\textbf{B}. $X=CA^{-1}$.}
\end{frame}

\begin{frame}{MC 6 - A $3\times3$ determinant}
Expand along the first row:
\begin{align*}
\det\begin{bmatrix}1&2&1\\2&1&3\\1&1&2\end{bmatrix}
&=1(2-3)-2(4-3)+(2-1)\\
&=-1-2+1=-2.
\end{align*}
\answerbox{\textbf{A}. $-2$.}
\end{frame}

\begin{frame}{MC 7 - Scaling and transpose}
For a $3\times3$ matrix,
\[
\det(2A^T)=2^3\det(A^T).
\]
Since $\det(A^T)=\det A=-3$,
\[
\det(2A^T)=8(-3)=-24.
\]
\answerbox{\textbf{D}. $-24$.}
\end{frame}

\begin{frame}{MC 8 - Membership in a span}
Let
\[
\vect v_1=\begin{bmatrix}1\\0\\2\end{bmatrix},
\qquad
\vect v_2=\begin{bmatrix}0\\1\\-1\end{bmatrix}.
\]
The first choice satisfies
\[
2\vect v_1-\vect v_2
=\begin{bmatrix}2\\-1\\5\end{bmatrix}.
\]
\answerbox{\textbf{A}. $\begin{bmatrix}2&-1&5\end{bmatrix}^T$.}
\end{frame}

\begin{frame}{MC 9 - Free variables}
A matrix has $7$ columns and $2$ pivots.
\[
\#\text{ free variables}
=\#\text{ columns}-\#\text{ pivots}
=7-2=5.
\]
\answerbox{\textbf{B}. Five free variables.}
\end{frame}

\begin{frame}{MC 10 - Block multiplication}
Use the second block row and the first block column:
\[
\begin{bmatrix}P&0\\Q&R\end{bmatrix}
\begin{bmatrix}S&0\\U&V\end{bmatrix}
=
\begin{bmatrix}
PS&0\\
QS+RU&RV
\end{bmatrix}.
\]
\answerbox{\textbf{D}. The lower-left block is $QS+RU$.}
\end{frame}

\subsection{Brief answers}
\sectioncard{Section II}{Brief Answers - 20 points}

\begin{frame}{Brief 1 - Rank and free columns}
For
\[
E=\begin{bmatrix}
1&0&2&0&-1\\
0&1&-3&0&4\\
0&0&0&1&2\\
0&0&0&0&0
\end{bmatrix},
\]
the pivot columns are $1,2,4$.
\answerbox{$\rank E=\boxed{3}$; free-variable columns $\boxed{3\text{ and }5}$.}
\end{frame}

\begin{frame}{Brief 2 - Standard matrix}
For
\[
T(x,y)=(x+2y,\ 3x-y),
\]
the columns of the standard matrix are
\[
T(1,0)=\begin{bmatrix}1\\3\end{bmatrix},
\qquad
T(0,1)=\begin{bmatrix}2\\-1\end{bmatrix}.
\]
\answerbox{$\displaystyle \boxed{[T]=\begin{bmatrix}1&2\\3&-1\end{bmatrix}}$.}
\end{frame}

\begin{frame}{Brief 3 - Solving a matrix equation}
Starting with
\[
A(X-I)=B,
\]
multiply on the left by $A^{-1}$:
\[
X-I=A^{-1}B.
\]
Add $I$ to both sides.
\answerbox{$\boxed{X=I+A^{-1}B}$.}
\end{frame}

\begin{frame}{Brief 4 - Determinant of a product}
Given $\det A=2$ and $\det B=-3$,
\begin{align*}
\det(A^{-1}B)
&=\det(A^{-1})\det B\\
&=\frac{1}{\det A}\det B\\
&=\frac12(-3).
\end{align*}
\answerbox{$\boxed{-\frac32}$.}
\end{frame}

\begin{frame}{Brief 5 - Product size}
If $A$ is $4\times3$ and $B$ is $3\times2$, then
\[
(4\times3)(3\times2)
\]
is defined because the inner dimensions agree. The outer dimensions give the product size.
\answerbox{$AB$ is $\boxed{4\times2}$.}
\end{frame}

\subsection{Exercises}
\sectioncard{Section III}{Exercises - 60 points}

\begin{frame}{Exercise 1 - Gaussian elimination}
\small
\[
\left[\begin{array}{ccc|c}
1&1&1&4\\2&-1&1&8\\-1&2&1&-1
\end{array}\right]
\xrightarrow[\ R_3\leftarrow R_3+R_1]{R_2\leftarrow R_2-2R_1}
\left[\begin{array}{ccc|c}
1&1&1&4\\0&-3&-1&0\\0&3&2&3
\end{array}\right].
\]
Use $R_3\leftarrow R_3+R_2$:
\[
\left[\begin{array}{ccc|c}
1&1&1&4\\0&-3&-1&0\\0&0&1&3
\end{array}\right].
\]
Then $z=3$, $y=-1$, and $x=2$.
\answerbox{$\boxed{(x,y,z)=(2,-1,3)}$.}
\end{frame}

\begin{frame}{Exercise 2 - Matrix expression}
Given
\[
A=\begin{bmatrix}0&1\\-2&3\end{bmatrix},
\qquad
B=\begin{bmatrix}1&-1\\4&0\end{bmatrix},
\]
first compute
\[
A^2=\begin{bmatrix}-2&3\\-6&7\end{bmatrix},
\qquad
AB=\begin{bmatrix}4&0\\10&2\end{bmatrix}.
\]
Therefore
\[
A^2+AB+I_2
=\boxed{\begin{bmatrix}3&3\\4&10\end{bmatrix}}.
\]
\end{frame}

\begin{frame}{Exercise 3 - $LU$ factorization}
The elimination multipliers are
\[
m_{21}=3,
\qquad
m_{31}=-2,
\qquad
m_{32}=2.
\]
Place them below the diagonal of $L$ and keep the echelon matrix as $U$.
\answerbox{$\displaystyle
\boxed{L=\begin{bmatrix}1&0&0\\3&1&0\\-2&2&1\end{bmatrix}},
\qquad
\boxed{U=\begin{bmatrix}1&2&-1\\0&4&3\\0&0&5\end{bmatrix}}$.}
\end{frame}

\begin{frame}{Exercise 4 - Determinant properties}
Since
\[
A=\begin{bmatrix}1&1&0\\0&2&1\\0&0&3\end{bmatrix}
\]
is triangular,
\[
\det A=1(2)(3)=6.
\]
For a $3\times3$ matrix,
\[
\det(-A)=(-1)^3\det A=\boxed{-6},
\]
\[
\det(3A)=3^3\det A=27(6)=\boxed{162}.
\]
\end{frame}

\begin{frame}{Exercise 4 - The matrix $C=3A+I_3$}
Compute $C$ directly:
\[
C=3A+I_3
=\begin{bmatrix}
4&3&0\\0&7&3\\0&0&10
\end{bmatrix}.
\]
The matrix is triangular, so
\[
\det C=4(7)(10)=\boxed{280}.
\]
\answerbox{$\det(-A)=-6$, $\det B=162$, and $\det C=280$.}
\end{frame}

\begin{frame}{Exercise 5 - Sparse determinant}
Expand along column $4$, then along the second row of the remaining determinant:
\begin{align*}
\det D
&=5\begin{vmatrix}1&0&2\\0&3&0\\4&0&-1\end{vmatrix}\\
&=5(3)\begin{vmatrix}1&2\\4&-1\end{vmatrix}\\
&=15(-1-8).
\end{align*}
\answerbox{$\boxed{\det D=-135}$.}
\end{frame}

\begin{frame}{Exercise 6(a) - Row echelon form}
\small
Starting with
\[
\left[\begin{array}{ccc|c}
1&1&1&2\\
2&3&4&5\\
3&4&a+5&b
\end{array}\right],
\]
apply
\[
R_2\leftarrow R_2-2R_1,
\quad
R_3\leftarrow R_3-3R_1,
\quad
R_3\leftarrow R_3-R_2.
\]
This gives
\[
\boxed{\left[\begin{array}{ccc|c}
1&1&1&2\\
0&1&2&1\\
0&0&a&b-7
\end{array}\right]}.
\]
\end{frame}

\begin{frame}{Exercise 6(b) - Classification}
The final row is
\[
\begin{bmatrix}0&0&a&b-7\end{bmatrix}.
\]
Therefore:
\begin{itemize}
  \item If $a\ne0$, the system has one solution.
  \item If $a=0$ and $b\ne7$, the equation $0=b-7$ is a contradiction, so there is no solution.
  \item If $a=0$ and $b=7$, the final row is zero and $z$ is free, so there are infinitely many solutions.
\end{itemize}
\end{frame}

\begin{frame}{Exercise 6(c) - The case $a=2$, $b=9$}
Substitute $a=2$ and $b=9$ into the echelon system:
\[
\left[\begin{array}{ccc|c}
1&1&1&2\\
0&1&2&1\\
0&0&2&2
\end{array}\right].
\]
Back substitution gives
\[
z=1,
\qquad y=-1,
\qquad x=2.
\]
\answerbox{$\boxed{(x,y,z)=(2,-1,1)}$.}
\end{frame}

\begin{frame}{Compact answer key}
\small
\textbf{Multiple choice}
\[
\boxed{\text{B, B, B, A, B, A, D, A, B, D}}
\]
\textbf{Brief answers}
\[
\rank E=3;\quad
[T]=\begin{bmatrix}1&2\\3&-1\end{bmatrix};\quad
X=I+A^{-1}B;\quad
-\frac32;\quad
4\times2.
\]
\textbf{Exercises}
\[
(2,-1,3),\quad
\begin{bmatrix}3&3\\4&10\end{bmatrix},\quad
A=LU,\quad
(-6,162,280),\quad
-135,
\]
with the parameter cases classified in Exercise 6.
\end{frame}

\begin{frame}{Final checks}
Before submitting a linear algebra solution:
\begin{itemize}
  \item state the row operations or determinant property used;
  \item keep matrix multiplication in the correct order;
  \item distinguish a singular coefficient matrix from an inconsistent system;
  \item write the final answer clearly and verify it when time permits.
\end{itemize}
\end{frame}

\end{document}
